\documentclass{article}
\usepackage[T1]{fontenc}
\usepackage{lmodern}
\usepackage{graphicx}
\usepackage{amsmath,amssymb}
\usepackage[a4paper,margin=2cm]{geometry}
\usepackage[absolute,overlay]{textpos}
\setlength{\TPHorizModule}{1mm}
\setlength{\TPVertModule}{1mm}
\usepackage[indonesian]{babel}
\usepackage{hyperref}
\setlength{\emergencystretch}{2em}
% unit: Halaman 16

\begin{document}

% === Halaman 16 (python/native) ===

• Mula-mula, semua byte dari kunci eksternal, K[0..b – 1], disalin ke dalam larik L yang 

berukuran c word, L[0.. c – 1] .  (ingat: b dalam byte, c dalam word, c < b)

• lalu padding dengan sejumlah 0 jika perlu (padding terjadi jika b bukan kelipatan w). 

• Kemudian inisialisasi larik S sebagai berikut:

 S[0]  Pw

 for i  1 to t – 1  do

     S[i]  S[i – 1] + Qw

 endfor

16


Parameter 

Simbol 

Nilai yang dibolehkan 

Ukuran blok (dalam bit) 

w 

16, 32, 64 

Jumlah putaran 

r 

0, 1, …, 255 

Panjang kunci eksternal K 

(dalam byte, 1 byte = 8 bit) 

b 

0, 1, …, 255

\end{document}
